\chapter*{Resumen}

%% Si se quiere enviar este resumen en una p\'agina a alguien, descomentar esto.

%\begin{center}
%\textbf{\large Universidad Argentina De la Empresa}

%\textbf{Facultad de Ingenier\'{\i}a y Ciencias Exactas}

%\textbf{Departamento de Tecnolog\'{\i}a Inform\'atica}

%\end{center}

%\begin{center}
%\large Ingenier\'{\i}a en Inform\'atica
%\end{center}

%\begin{center}
%\textbf{\large Título de la PFI}
%\end{center}

%\begin{center}
%Apellido, Nombres
%Apellido, Nombres
%\end{center}

%\noindent

El presente trabajo aborda la circulación de desinformación en las redes sociales y la falta de una herramienta gratuita, adaptada al contexto argentino, que permita evaluar la veracidad de una publicación en el momento en que se la lee. Para ello, se desarrolló Factum, una extensión de Google Chrome para Twitter/X que combina un clasificador de lenguaje natural ajustado sobre corpus en español, las señales públicas de la cuenta autora y el contraste de la afirmación con fuentes oficiales, medios de referencia y verificadores profesionales.

Cada resultado se acompaña de una justificación y del enlace a las fuentes que lo sostienen, de modo que el usuario pueda comprobarlo por su cuenta. La solución se implementó como un prototipo funcional que se integra a la línea de tiempo de Twitter/X y se complementa con un panel web, que reúne el histórico personal de análisis y ofrece a las organizaciones una vista agregada de las tendencias. Sobre 108 publicaciones argentinas verificadas, el veredicto del sistema alcanzó un valor F1 macro de 0,654, superior al de la línea base (0,527) y por debajo de la meta de 0,75; el clasificador por sí solo obtuvo 0,540, lo que muestra la distancia entre las noticias de los corpus académicos y las publicaciones breves de la red social y explica que el veredicto dependa del contraste con fuentes. El análisis a demanda demora 14,9~s en el percentil 95, por encima de los ocho segundos previstos. La prueba de usabilidad, con tres participantes, obtuvo un puntaje SUS medio de 74,2.